\documentclass[aspectratio=169]{beamer}

\usetheme{Madrid}
\usecolortheme{seahorse}

\usepackage{graphicx}
\usepackage{xcolor}
\usepackage{hyperref}

\title{Meeting 1: Welcome to Arcturus}
\subtitle{Autonomy Onboarding}
\author{MIT Arcturus}
\date{}

\begin{document}

\begin{frame}
    \titlepage
\end{frame}

% ============================================================
\section{The Boat}
% ============================================================

\begin{frame}{What We Do}
    We build an autonomous boat that competes at RoboBoat every year.
    \vfill
    % TODO: photo/video of the boat
    \begin{center}
        \textit{(boat photo or competition footage here)}
    \end{center}
\end{frame}

\begin{frame}{What the Boat Has to Do}
    The competition has several tasks:
    \vfill
    \begin{itemize}
        \item Follow the path --- drive through pairs of red and green buoys
        \item Speed challenge --- same thing, but fast
        \item Docking --- park in a specific spot
        \item Delivery --- drop something off at a target
    \end{itemize}
    \vfill
    % TODO: competition video or task diagram
    Today we'll work on a simplified version of ``follow the path.''
\end{frame}

\begin{frame}{How the Boat Works (Big Picture)}
    \begin{center}
        \Large
        Perception $\rightarrow$ Planning $\rightarrow$ Control
    \end{center}
    \vfill
    \begin{enumerate}
        \item Perception: cameras and LiDAR detect what's around the boat
        \item Planning: software figures out where to go next
        \item Control: motors drive the boat to that location
    \end{enumerate}
    \vfill
    All of this is software that you will be writing and improving.
\end{frame}

% ============================================================
\section{The Simulation}
% ============================================================

\begin{frame}{Let's Try It}
    Before we explain anything, let's just run the simulation.
    \vfill
    \begin{enumerate}
        \item Open your Docker container
        \item Build the workspace and launch the sim
        \item Drive the boat around with the keyboard (W/Q/E)
    \end{enumerate}
    \vfill
    Follow the setup instructions in the lab doc.
    \vfill
    % TODO: screenshot of the sim in RViz
\end{frame}

\begin{frame}{What Just Happened?}
    When you launched the sim, several programs started running:
    \vfill
    \begin{itemize}
        \item One program simulates the boat's physics
        \item One program reads your keyboard input
        \item One program shows everything in RViz
    \end{itemize}
    \vfill
    These programs are running at the same time and sending messages to each other.
    \vfill
    Let's look at those messages. Try running:
    \begin{center}
        \texttt{ros2 topic list}
    \end{center}
\end{frame}

% ============================================================
\section{ROS 2}
% ============================================================

\begin{frame}{What You Just Saw}
    \texttt{ros2 topic list} showed you all the message channels in the system.
    \vfill
    Try this:
    \begin{center}
        \texttt{ros2 topic echo /boat\_pose}
    \end{center}
    \vfill
    Drive the boat around while this is running. You'll see the position values change in real time.
    \vfill
    These message channels are called \textit{topics}.
\end{frame}

\begin{frame}{Nodes and Topics}
    Now we can put names to what we saw:
    \vfill
    \begin{itemize}
        \item Each program is called a \textit{node}
        \item Nodes send messages to each other through \textit{topics}
        \item Sending a message is called \textit{publishing}
        \item Receiving messages is called \textit{subscribing}
    \end{itemize}
    \vfill
    Example from our sim:
    \begin{center}
        Keyboard $\xrightarrow{\text{velocity commands}}$ Simulation $\xrightarrow{\text{boat position}}$ Visualization
    \end{center}
\end{frame}

\begin{frame}{Why This Matters}
    The keyboard program sends velocity commands on a topic.
    \vfill
    The simulation doesn't know or care where those commands come from --- it just listens to the topic.
    \vfill
    So we can write our own program that sends commands to the same topic, and the simulation will follow them.
    \vfill
    That's how we go from keyboard control to autonomous control.
\end{frame}

% ============================================================
\section{Hands-On}
% ============================================================

\begin{frame}{Your First Node}
    Goal: write a program that moves the boat forward automatically.
    \vfill
    \begin{enumerate}
        \item Create a new Python file
        \item Make it send a velocity command every 0.5 seconds
        \item Register it so ROS 2 can find it
        \item Run it and watch the boat move on its own
    \end{enumerate}
    \vfill
    Follow the lab doc for the code and step-by-step instructions.
\end{frame}

\begin{frame}{Waypoint Following}
    Moving forward is a start. Let's make the boat go to a specific point.
    \vfill
    The idea:
    \begin{enumerate}
        \item Know where the boat is (read its position from a topic)
        \item Know where it should go (click a point in the visualizer)
        \item If we're far from the target, speed up
        \item If we're close, slow down
        \item If we're pointed the wrong way, turn
    \end{enumerate}
    \vfill
    This is called PID control --- steer harder when you're far, ease off when you're close.
\end{frame}

\begin{frame}{PID Control}
    \begin{center}
        \Large
        correction = (how far off) + (for how long) + (how fast approaching)
    \end{center}
    \vfill
    \begin{itemize}
        \item P --- far from target? Steer hard. Close? Steer gently.
        \item I --- been slightly off for a while? Nudge a bit more.
        \item D --- approaching fast? Ease off so you don't overshoot.
    \end{itemize}
    \vfill
    We give you a PID class. You feed it the error, it gives you back a motor command.
\end{frame}

\begin{frame}{Navigating Through Buoys}
    Final challenge: drive through a course of red and green buoy pairs.
    \vfill
    \begin{enumerate}
        \item Find the next pair of buoys ahead of the boat
        \item Compute the midpoint between the red and green buoy
        \item Use your waypoint follower to drive there
        \item When you arrive, move on to the next pair
        \item When there are no more pairs, stop
    \end{enumerate}
    \vfill
    This is a simplified version of what the real boat does at competition.
\end{frame}

\begin{frame}{What You Built Today}
    \begin{center}
        \Large
        Detect buoy pair $\rightarrow$ Find midpoint $\rightarrow$ PID to midpoint $\rightarrow$ Repeat
    \end{center}
    \vfill
    On the real boat, ``detect buoy pair'' uses cameras and LiDAR instead of known positions, and ``PID to midpoint'' uses path planning to avoid obstacles.
    \vfill
    Next meeting: we'll look at how the real system works.
\end{frame}

\end{document}
